\chapter{NHẬN XÉT VÀ ĐỀ XUẤT}

\section{Ưu điểm}

Đề tài bám đúng chuyên ngành Kỹ thuật Điện tử Viễn thông: đọc mã linh kiện, đối
chiếu datasheet và xem thiết bị tự động hóa được chọn như thế nào trong một
doanh nghiệp cung ứng thật. Danh mục công khai của TECHPAL đủ rộng để em không
phải dừng ở mô tả chung, mà đi vào bốn nhóm sản phẩm và bốn mã đại diện.

Việc tra cứu MAX487CSA+T, BMXAMI0410, 6ES7153-1AA03-0XB0 và CR1000X giúp em
thấy cùng một bài toán chức năng có thể đi theo các hướng khác nhau. IC rời
phù hợp thiết kế nhúng; mô-đun analog phải khớp rack và CPU; bộ ghi dữ liệu
tích hợp tiện cho trạm đo độc lập. Trường hợp Siemens phân loại
6ES7153-1AA03-0XB0 là phụ tùng là ví dụ rõ nhất: nếu chỉ nhìn tên chức năng
thì dễ đề xuất sai cho dự án mới.

Em cũng tập được thói quen không tin mô tả ngắn trên trang bán hàng. Bảng kiểm
kỹ thuật (Bảng~\ref{tab:technical-checklist}) buộc phải trả lời đủ định danh,
thông số điện, giao tiếp, cơ khí, môi trường, vòng đời và chứng từ trước khi
chốt mã. Bài toán thu thập nhiệt độ tủ kỹ thuật dù chỉ ở mức phương án giấy
vẫn hữu ích, vì nó buộc em đi từ yêu cầu chức năng sang kiến trúc rồi mới so
sánh vai trò từng mã. Các kỹ năng đọc tiếng Anh chuyên ngành, ghi chép nguồn
và trình bày phương án là phần em thấy mang sang các học phần thiết kế sau này
được.

\section{Nhược điểm}

Về điều kiện thực tập, thời gian ngắn nên em chỉ kịp đi sâu một mảng cung ứng.
Em không tiếp cận được quy trình nội bộ, dữ liệu kho hay phòng thí nghiệm, vì
vậy Hình~\ref{fig:supply-flow} chỉ là mô hình hóa từ thông tin công khai, không
phải luồng đang áp dụng tại công ty. Bài toán nhiệt độ cũng chưa chọn cảm biến,
chưa có số liệu đo, chưa đánh giá được sai số, nhiễu hay độ bền khi vận hành.

Về phía bản thân, danh mục công nghiệp lớn và cách đặt mã khác nhau giữa các
hãng khiến em dễ chú ý vài thông số nổi (số kênh, điện áp) rồi bỏ sót chi tiết
quyết định. Hậu tố ``T'' của MAX487CSA+T liên quan quy cách đóng gói tape and
reel, không phải biến thể mạch; nếu đặt hàng thiếu hậu tố có thể nhận sai dạng
cuộn. Với BMXAMI0410, em chưa đọc được tài liệu tương thích rack/CPU ở mức
triển khai. Với mã Siemens, nếu không tra trạng thái vòng đời thì sẽ không biết
sản phẩm chỉ còn dạng phụ tùng. Do đó mọi kết luận về khả năng thay thế trong
báo cáo đều cần người có kinh nghiệm xem lại.

Nhược điểm này cũng thể hiện ở cách viết: một số đoạn còn kể lại thông tin web
nhiều hơn là phân tích quyết định kỹ thuật. Em nhận ra rằng khảo sát tài liệu
là bước cần, nhưng chưa đủ để khẳng định một mã ``dùng được'' trong hệ thống
cụ thể.

\section{Đề xuất}

Với đơn vị tiếp nhận, em đề xuất có sẵn một bộ tài liệu ngắn cho sinh viên mới:
quy tắc đọc mã và hậu tố, mẫu so sánh thiết bị theo các cột của bảng kiểm, hướng
dẫn chống tĩnh điện (ESD), kèm một tình huống chọn mã thật -- ví dụ một yêu cầu
thay thế hoặc một đơn hàng đã xử lý -- để sinh viên quan sát cách Sale, Mua hàng
Quốc tế và Logistics giữ thông tin nhất quán. Nếu điều kiện cho phép, cho sinh
viên xem một lần kiểm tra ngoại quan hoặc đối chiếu chứng từ khi nhận hàng sẽ
giúp nối phần đọc datasheet với công việc thực tế, dù không cần giao việc nội
bộ.

Với bản thân, em sẽ hoàn thiện bảng kiểm cho đủ bốn mã đã chọn, đặc biệt các
cột vòng đời và tương thích hệ thống còn mỏng. Nếu có điều kiện, em muốn làm
một nguyên mẫu nhỏ đo nhiệt độ tủ kỹ thuật để kiểm chứng phương án giấy: chọn
cảm biến, xác định dải đo, chu kỳ lưu và cổng truyền về máy tính. Song song đó,
em cần luyện viết ghi chú kỹ thuật ngắn hơn, tập trung vào tiêu chí chốt mã thay
vì thuật lại toàn bộ trang sản phẩm. Những việc này sẽ giúp em chắc hơn khi học
tiếp thiết kế mạch, hệ thống nhúng và tự động hóa.
